% Related Work. Grounded in docs/positioning.md; citations in related_work.bib.
% Draft prose --- tighten to ~700 words at revision.

\paragraph{Sampling-based hallucination detection.}
Semantic entropy estimates a model's uncertainty by sampling several answers,
clustering them into meaning-equivalence classes, and taking an entropy over those
classes~\citep{kuhn2023semantic,farquhar2024detecting}. How the classes are \emph{weighted}
is not a detail, and the originating papers do not settle on one answer:
\citet{kuhn2023semantic} average cluster surprisals under length-normalised sequence
likelihoods, \citet{farquhar2024detecting} normalise those likelihoods into a categorical
distribution and take its plug-in entropy (their Eq.~(5), which they designate as their main
method), and the same paper introduces a third, \emph{discrete} variant that weights each
cluster by the number of generations in it. The three do not share their measurement
properties, so we name ours: we study the discrete variant, for reasons and with consequences
set out in Section~\ref{sec:methods}. Its central
selling point is \emph{linguistic invariance}: because clustering is on meaning
rather than surface form, the score is meant to be stable across rephrasings.
Subsequent work improves the estimator. Semantic Reformulation Entropy (SRE)
adds input-side reformulations and energy-based hybrid clustering for
robustness~\citep{tong2025sre}, and Semantic Entropy Probes (SEPs) approximate
the entropy from a single hidden state, removing the sampling
cost~\citep{kossen2024seps}. A parallel line questions the \emph{clustering} step
itself: Kernel Language Entropy replaces hard NLI equivalence-class clustering with a
semantic-graph kernel, precisely because hard-NLI clustering is
brittle~\citep{nikitin2024kle} --- the very sensitivity our attack exploits. We treat
these as \emph{detectors to be attacked}: our contribution is to test the invariance claim
adversarially, a question none of these papers ask of themselves. Whether the attack
transfers to such richer-clustering variants, or is specific to hard-NLI-equivalence
clustering, is an open question our null-controlled protocol is designed to pose.

\paragraph{Hallucination-elicitation attacks on models.}
The closest methodological neighbours attack the \emph{generating model}, not a
detector. SECA frames realistic hallucination elicitation as constrained
optimization over the prompt under semantic-equivalence and coherence
constraints, solved with a zeroth-order method~\citep{liang2025seca}; its
successor REALISTA extends this to free-form reasoning models via continuous
latent edit directions~\citep{liang2026realista}. We reuse SECA's optimization
loop as tooling, but our objective is the \emph{detector's} uncertainty score
rather than the model's answer correctness, and we add a direction (false alarm)
absent from both. Crucially, neither evaluates a hallucination detector, so
neither subsumes our result despite the shared machinery.

\paragraph{Paraphrase attacks on text-origin detectors.}
A mature line paraphrases text to evade AI-versus-human
classifiers~\citep{cheng2025advparaphrasing,fang2025copa}. \citet{wang2026depo} give the
cleanest formalisation of the primitive, casting detector-evasive paraphrasing as a
constrained Markov decision process in which semantic preservation is an explicit
constraint rather than a scalarised reward term. These share the
high-level move ``paraphrase to fool a detector'' but differ on every axis that
matters here: the victim is a text-origin classifier, not a sampling-based
uncertainty estimator; the attacks are evasion-only (one-directional); and the
equivalence guarantee is a soft similarity score rather than the hard
bidirectional-NLI gate we require. We adopt a distinct name to avoid confusion
with~\citet{cheng2025advparaphrasing}.

\paragraph{Manipulating uncertainty and probes.}
The closest input-side neighbour is \citet{obadinma2025verbal}, who show that
semantics-preserving perturbations degrade an LLM's \emph{verbal} confidence---carrying that
group's own calibration-attack line~\citep{obadinma2024calibration} from vision into language. The victim
differs in the way that matters here: verbal confidence is the model's prompted
self-report, whereas semantic entropy is computed from a sampled answer set clustered by
meaning, so the attack surface is the sampling-and-clustering pipeline rather than a
self-assessment; we also impose a bidirectional-NLI equivalence gate they do not, and our
success criterion \emph{requires} the model's answer status to survive the paraphrase,
whereas their attacks frequently change the answer. Also on the ``manipulate uncertainty''
axis, \citet{uncertaintyfragile2024}
shift answer confidence via a fine-tuned backdoor trigger on multiple-choice
self-evaluation --- a different threat model (weight access, not input
paraphrase).

We should be explicit that the \emph{direction} we foreground is not itself new. Moving a
confidence signal in an attacker-chosen direction while holding the underlying prediction
fixed is an established family with a name: \citet{obadinma2024calibration} give what they
describe as the first dedicated study of \emph{calibration attacks}, which ``trap victim
models to be heavily miscalibrated without altering their predicted labels'', and the family
is still active---\citet{martinezmartinez2026uat} propose an underconfidence attack that
lowers confidence while preserving the prediction, paired with an adversarial-training
defence. Forcing a detector into a false \emph{positive} has a
direct text-modality precedent: \citet{rusert2025redherring} perturb text so that an
attack-detection model flags an input while the classifier's prediction stays correct.
Concurrently, \citet{khanmohammadi2026answerpreserving} manipulate deployed confidence
channels bidirectionally under an answer-preservation constraint, including an inflation
direction, and report the same qualitative collapse of a gated system toward its ungated
floor. We therefore do not claim the false-alarm direction as a discovery. What we claim is
its instantiation against a \emph{sampling-based} detector, which differs on two axes that
carry consequences rather than labels. First, the manipulated quantity is an entropy over
meaning-clustered samples rather than a softmax or a learned detector's logit: a stochastic
finite-sample estimate, bounded above by $\log N$, produced by the same NLI model that
certifies the attacker's constraint. Second, the perturbation is a question rewrite rather
than a pixel budget or a hidden-state intervention---channels a deployed adversary does not
possess---so semantic invariance cannot be granted by an $\ell_p$ ball and must instead be
certified. We accordingly treat the equivalence constraint as an instrument requiring its
own validation rather than as a definition: we constrain, re-check that answer invariance
actually held, measure how often the gate nonetheless admitted a non-equivalent rewrite and
report that as a disclosed lower bound, and validate the oracle making the judgement on an
adversarial discrimination that disqualified the obvious cheap candidate ($0.51$). An attack
whose constraint is an $\ell_p$ ball has no analogue of these steps, because there the
constraint is definitional rather than inferred. Most relevant of all, CORVUS red-teams single-pass hallucination
detectors and degrades SEP~\citep{corvus2026} --- establishing that SE-\emph{family}
detectors can be attacked, which we do not claim as our novelty. Our distinction is
substantive rather than a label: CORVUS perturbs a frozen hidden-state probe
\emph{model-side} (LoRA fine-tuning, fixed inputs, suppression-only, no
semantic-equivalence constraint, Alpaca/FAVA), whereas we perturb the \emph{input} under a
hard equivalence constraint and target the sampling-and-NLI-clustering pipeline itself --- a
different attack surface --- in both directions, on TriviaQA/SQuAD. A weaker neighbour prompts
LLMs to fool generic activation probes~\citep{blandfort2025redteamprobes}, but
not in the hallucination-detection setting. Each ingredient of our attack is by now standard: constrained search over meaning-preserving
paraphrases is an established primitive---\citet{kaneko2026paa} run essentially our optimiser
loop, searching paraphrases that raise an LLM reviewer's score under semantic-equivalence and
naturalness constraints---and directional manipulation of a confidence signal
under prediction-preservation is an established family. What does not carry over is the
\emph{evaluation}: a victim whose score is a stochastic, ceiling-bounded estimate produced by
the same model that certifies the attacker's constraint invalidates the standard protocol in
three specific ways --- selection on estimator noise, right-censoring at the cap, and
unattributability under a shared clusterer --- none of which arises for a classifier victim.
That is where we locate our contribution (we evaluate the false-alarm direction here and
build the hide direction as its mirror). The bidirectional-NLI equivalence gate is our design choice for
certifying meaning-preservation, not itself a novelty axis --- indeed our own analysis
treats the detector's reliance on that same NLI model as the central confound.

\paragraph{What the estimator can and cannot express.}
A separate line studies semantic entropy as a \emph{statistic} rather than as a detector.
\citet{mccabe2025alphabet} show that the discrete plug-in estimator underestimates the true
semantic entropy at practical sample sizes, because the semantic alphabet is under-covered,
and correct it upward; \citet{pan2026shade} pursue the same correction with a
coverage-adjusted estimator fusing Good--Turing coverage with a spectral term over an
entailment-weighted response graph. We characterise the same estimator's opposite tail---it
cannot exceed $\log N$, and in deployment it piles up there. That two groups now correct the
estimator upward is itself evidence for our reading rather than against it: a remedy at the
level of the \emph{estimator} concedes that the deployed plug-in statistic is the limiting
factor, which is precisely our claim. The two accounts are complementary rather than
competing, and we are careful not to promote that into an explanation. Two mechanisms can
make a plug-in estimate fall short of the semantic entropy it targets: truncation at
$\log N$, which binds whenever the true quantity exceeds the cap, and the finite-sample
negative bias of a plug-in entropy over an alphabet the sample under-covers, which is the
mechanism these papers propose and correct for. They are not mutually exclusive, both are
active at deployed $N$, and nothing we measure apportions the understatement between them,
so we claim agreement with their finding and not a diagnosis of it.
Separately, \citet{sun2026granularity} name the \emph{score granularity
gap} for black-box LLM classification: a confidence score may rank well and still take so
few distinct values that an operator has only a handful of usable thresholds. \emph{Discrete}
semantic entropy inherits this pathology at the budget it is deployed at---$39$ attainable
values at the standard $N{=}10$, two of them in the top tenth of the range---with the twist
that the granularity limit is set by the sampling budget rather than by a verbalisation
format, and is therefore relaxable by spending in a way a verbalisation format is not:
$N{=}20$ affords $455$ values, seven of them in that top tenth. We accordingly assert the
claim at $N{=}10$ and nowhere else, and what carries to larger $N$ is only that the lattice
stays sparsest at the top of the scale, which is where a false alarm and an operator's
threshold both have to land. Our contribution here is the instantiation and its consequence
for attack evaluation, not the argument itself. The
likelihood-weighted estimator does not inherit it: real-valued cluster weights are not
confined to the entropies of integer partitions of $N$, so the lattice dissolves even though
the $\log N$ ceiling remains (Section~\ref{sec:methods}). The granularity claim is thus the
sharper and the narrower of our two---narrower in the estimator it holds for and in the
sample budget it is asserted at---and we keep the two apart wherever we make them. Read
together with the originating papers---which report that clusters are few and grow slowly
with $N$, and treat that as a computational
saving~\citep{kuhn2023semantic, farquhar2024detecting}---the picture is of a well-known
observation whose measurement consequences have not been drawn out.

\paragraph{Evaluating the evaluation.}
Because our framing is protocol-led, two methodological strands matter beyond the attack
neighbours. First, the individual failure modes our controls target are classical ---
selection on the evaluated system's own scores, selection of an extremum of a noisy
estimate, and success criteria that drift from the construct --- and so are the individual
remedies: comparing an adversarial perturbation against a random one of equivalent
magnitude is established practice in vision robustness
evaluation~\citep{rossolini2026worstcase}, and plain-paraphrase baselines are routine in
AI-text-detector evasion~\citep{cheng2025advparaphrasing}. We claim neither as new. That an
uncontrolled baseline can manufacture an entire effect is not hypothetical:
\citet{zheng2026matchedctrl} report an apparent gain of up to $31.8$ points from a
perturbation-based selection rule that disappears once a format- and budget-matched control
is introduced, the control tracking or exceeding the method on every benchmark. What we
are not aware of, for a paraphrase attack on a sampling-based hallucination detector, is
the \emph{joint} application of score-independent selection, answer-invariance checking,
and a noise floor that is \emph{budget-matched} to the attacker's own search --- the
quantity the attack's maximum is actually compared against. Our contribution is assembling
these into a pre-registered protocol with a decision rule, and we state the gap as ``to our
knowledge'' rather than as priority.
Second, using an LLM as an equivalence oracle inherits the LLM-as-judge literature's
lesson that judges track humans well yet carry biases and must be validated per
task~\citep{zheng2023judging}. Our judge contribution is accordingly not the concept but
the \emph{adversarial-case} validation: the oracle is accepted only on the hard-negative
discrimination where the cheap embedding oracle demonstrably fails ($0.93$ vs.\ $0.51$),
with its pre-registered acceptance gate, the failed positive-recognition leg, and the
consequent false-alarm-only re-scoping disclosed rather than assumed away (Methods).
